\scanpage{481}

Pour le deuxième point, on a une réponse immédiate (supposant connus déjà les
$\unsure{\mathcal{N}_{g,\nu}|N_{g,\nu}}$). Soit \unsure{une|un} $(\pi,a)$ de type $g,\nu$,
avec \deleted{\unsure{quand}} relèvement \inserted{partiel} de $\Pi_a$ \unsure{en|un}
$\Gamma'_a\hookrightarrow \unsure{\mathcal{N}_a|N_a}$.
Elle sera admissible \unsure{ssi} on peut trouver
un $(\pi',a')$ de type $(g,\nu')$, \inserted{et} un
plongement de $\pi'$ comme \unsure{ss-groupe|ss groupe} d'indice fini de $\widehat{\pi}_{03}$ compatible avec $a'$,
\unsure{et} un sous-groupe $\unsure{\Gamma}$ ouvert \inserted{$\Gamma''$} de $\Pi_{03}$, et
une \unsure{action} effective \inserted{admissible} de $\Gamma''$ \deleted{\unsure{admissible}}
sur $\widehat{\pi}_{03}$ qui
relève l'action ext. donnée, et qui
\unsure{laisse} $\pi'$, de façon que, à isom
près, $(\pi,a)$ \inserted{et le germe d'action de $\Gamma'$ dessus}
se déduise de $(\pi',a')$
par l'action de $\Gamma''$ dessus, par ``bouchage''
d'un paquet de trous \inserted{(et \unsure{oubli} d'une action effective au profit
de l'action \unsure{extérieure})}. Il faut
dans cette approche simplement savoir
préciser \inserted{\unsure{algébriquement}} ce qu'on entend par
action ``admissible'' de \inserted{$\Gamma''\subset$} $\Pi_{0,3}$ sur $\widehat{\pi}_{0,3}$
--- sous-entendu que ça doit correspondre
aux pts de $\unsure{U_{0,3}}$ rat. sur \deleted{$\unsure{\mathbb{Q}}$} le corps
de nombres défini par $\Gamma''$. On peut conjecturer

\margin{480}

\scanpage{482}
\mspage{244}

celle-ci, par ``bouchage des trous''.

Il semble qu'on puisse sur le $\overline{\mathrm m}$
principe donner une description des
$\unsure{\mathcal{N}_{g,\nu}|N_{g,\nu}}$ (donc des $\unsure{\Pi_{g,\nu}|\Gamma_{g,\nu}}$)
en termes de
$\Pi_{0,3}$. Utilisant les hom de transition
$\unsure{\widehat{\widehat{\mathcal{T}}}_{g,\nu}}
\unsure{\twoheadrightarrow|\longrightarrow}
\widehat{\widehat{\mathcal{T}}}_{g,\nu'}$,
il suffit de le
faire, quand $g$ est fixé, pour des $\nu$
grands --- assez grands pour qu'il
existe une courbe \inserted{alg.} de type $g,\nu$ \inserted{sur $\mathbb{Q}$}
qui soit un rev. étale de $U_{0,3}$ sur $\mathbb{Q}$.

On considère donc un plongement \unsure{correspondant}
de $\pi_{g,\nu}$ dans $\pi_{0,3}$, et on
\unsure{décide|décrète} que si \deleted{$\unsure{\mathcal{T}_{03}}$} on peut faire
opérer $\mathrm{ext}^{t}$ $\Pi_{0,3}$ dans $\pi_{g,\nu}$, de
façon que $\pi_{g,\nu}\longrightarrow\pi_{0,3}$ \unsure{commute}
à ces actions ext, \uline{alors} le sous
groupe de $\widehat{\widehat{\mathcal{T}}}_{g,\nu}$ engendré par
$\widehat{\mathcal{T}}_{g,\nu}$ et
$\Pi_{0,3}$ est \deleted{$\unsure{\widehat{\mathcal{T}}}$}
$\unsure{\mathcal{N}_{g,\nu}|N_{g,\nu}}$.

\margin{481}

\scanpage{483}
\mspage{245}
34. Description heuristique profinie de la catégorie
des courbes alg. \unsure{définies} sur \deleted{\unsure{$\mathbb Q$}} des sous-
extensions finies $K$ de $\overline{\mathbb Q}_0/\mathbb Q$ (i.e. de $\mathbb C/\mathbb Q$).

On se borne aux courbes \inserted{géom.} connexes \inserted{(par commodité)}, anabéliennes
(par nécessité provisoire \unsure{!}, cf plus bas sur \unsure{la façon} de
se débarrasser de cette \unsure{condition|restriction}). La donnée d'
un rev. \unsure{universel} de \deleted{\illegible{1}} $U_{\overline{\mathbb Q}_0}$ définit une
structure d'extension
\[
(1)\qquad 1\longrightarrow\pi\longrightarrow\mathcal E\longrightarrow\unsure{\Gamma'|\Gamma^{+}}\longrightarrow1\qquad\text{ou}
\]
$\unsure{(2)|(1')}$ ou encore on a un hom $\mathcal E\longrightarrow\Pi$
(\unsure{image} $\unsure{\Gamma'|\Gamma^{+}}$ ouverte dans $\Pi$,
noyau \unsure{repéré} $\pi$),
où $\unsure{\Gamma'|\Gamma^{+}}\subset\Pi_{\mathbb Q}\simeq\Pi_{0,3}$ est le sous-gr. ouvert
correspondant à $K$. Se donner une telle
extension (moyennant $\unsure{\mathrm{Centre}}(\pi)=1$) revient à
se donner une action extérieure de $\Gamma'$ sur $\pi$.
Une première question : faut-il mettre la
structure à lacets de $\pi$ dans les \deleted{\illegible{1}} données \unsure{$\mathcal E$} de
l'objet (1) (ou \unsure{(1')}) censé décrire $U\unsure{\mathcal E}/K$ ?

\uline{Conjecture} Ce n'est pas la peine -- la structure à lacets
de $\pi$ est la seule structure à lacets \unsure{invariante}
par l'action extérieure de $\unsure{\Gamma'|\Gamma}$ (ou de $\unsure{\Gamma^{\flat}|\Gamma^{\natural}}$). Les
sous-groupes à lacets $L_i$ sont les sous-groupes
maximaux\inserted{dans $\pi$}, tels que $\mathrm{Norm}_{\mathcal E}(L_i)\longrightarrow\unsure{\Gamma'|\Gamma}$ ait
comme image un ss-groupe ouvert de $\unsure{\Gamma'|\Gamma}$.
\margin{\unsure{et} tout sous-groupe $\ne1$ de $\pi$ qui a cette propriété est contenu dans un unique $L_i\ldots$}
Je
présume aussi que tout hom \deleted{\unsure{d'}}\inserted{entre} extensions
$\mathcal E$ de $\unsure{\Gamma'|\Gamma}$ par un $\pi$, $\mathcal E'$ de $\unsure{\Gamma'|\Gamma}$ par un $\pi'$
($\mathcal E,\mathcal E'$ provenant de courbes algébriques $U,U'$) et tel
que l'image de $\mathcal E$ dans $\mathcal E'$ soit ouverte, respecte
\margin{482}

\scanpage{484}
nécessairement la structure à lacets, et en fait
provient d'un hom (unique, on le sait) de courbes
algébriques.

En tous cas, si on admet la description des
ss-groupes lacets, il \unsure{sera} clair que l'image par
$\unsure{u}$ d'un $L_i$ sera, ou bien $(1)$, ou bien contenue
dans un unique $L'_j\ldots$

Ceci signifierait que le foncteur des courbes
algébriques sur $K$\inserted{(morphismes les \unsure{applications dominantes})}, vers les groupes profinis ext.
sur lesquels $\Gamma$ opère, serait pleinement fidèle.
Le foncteur ``extension du corps de base'' de
$K$ à $K'$ correspond au foncteur : restriction
d'une op. extérieure (ou d'une structure
d'extension) de $\Gamma$ à $\Gamma'$ (ss-groupe ouvert).
Les rev. \inserted{\unsure{étales}} finis de $U$ correspondent aux
$\mathcal E$-ensembles : [$\widetilde U$ \unsure{étant choisi}]\ldots

Pour décrire l'image essentielle de ce
foncteur, on n'est pas réduit aux conjectures :
on part de l'extension \inserted{can.} $\mathcal E_{03}$ de $\Pi_{03}=\Pi$
par $\widehat{\pi}_{03}$, on prend un \unsure{sous-groupe|ss-groupe} \inserted{ouvert} $\mathcal E'$
de $\mathcal E_{03}$, d'image $\Gamma\subset\Pi_{03}$, noyau
$\pi'\subset\widehat{\pi}_{0,3}$, et dans la structure à lacets
canonique de $\pi'$\inserted{de genre $g$} (induite par celle de $\widehat{\pi}_{03}$)
on prend un $\unsure{I|\overline I}\subset I(\pi')$\inserted{tel que $2g+\unsure{\mathrm{Card}(I)}\geq3$}, \inserted{stable par $\Gamma$, et} on ``bouche les trous''
en $I(\pi')\setminus\unsure{I|\overline I}$. De même, pour décrire quand
\margin{483}

\scanpage{485}
\mspage{\unsure{246|245}}
une action effective, relevant une action extérieure
admissible, est elle-même admissible -- \deleted{\unsure{peut}}
i.e. peut être obtenue à partir d'une \uline{\unsure{courbe}}
algébrique $U$, \uline{ponctuée} par un pt rationnel sur $K$.
Ceci ne signifie pas pour autant, que si
$U$ est donnée par une action ext. de $\Gamma$ sur un $\pi$,
que les classes de conj. de relèvements admissibles
de cette action \unsure{proviennent} bien toutes de
points de $U$ rat/$K$. Mais on voit de suite que
ceci sera le cas, dès que l'on admet la
pleine fidélité pour les \uline{\unsure{monomorphismes}}

\uline{NB} Même pour les automorphismes de $U_{0,3}$
$=\mathbb P^1_{\mathbb Q}\setminus\{0,1,\infty\}$, cette pleine fidélité n'est pas
du tout claire. Il faudrait prouver que \deleted{\illegible{1}}
le commutant de $\Pi_{0,3}$ dans $\widehat{\widehat{\mathcal T}}_{0,3}$ est réduit
à $\unsure{\mathfrak S_3|\mathcal S_3}$. La situation est \unsure{moins sans espoir},
quand on se donne, avec la structure de
groupe \inserted{profini} extérieur à opérateurs $\Gamma$, une \unsure{authentisation}
de $\pi$ (ce qui suppose qu'on a \unsure{explicité}
une structure à lacets) invariante par \deleted{\illegible{1}} $\Gamma$
-- \unsure{donc} dans $\widehat{\widehat{\mathcal T}}(\pi)$ on dispose d'un $\unsure{N(\pi)|\mathcal N(\pi)}$,
et $\Gamma\hookrightarrow\unsure{N(\pi)|\mathcal N(\pi)}$. Dans ce point de vue, \inserted{pour} un
hom \inserted{$\pi'\longrightarrow\pi$} de $\Gamma$-groupes extérieurs (\unsure{authentisés}),
\unsure{il} est sous-entendu qu'il est compatible avec
les \unsure{authentisations}, i.e. si $\pi''$ est l'image de $\pi'$ dans $\pi$,
\margin{484}

\scanpage{486}

on voit que l'\unsure{authentisation} de $\pi''$ \unsure{déduite} de celle de $\pi'$ par ``passage au quotient'', soit celle \unsure{induite} par $\pi$. \deleted{\unsure{L'intérêt}} En particulier, \inserted{pour} les autom.\inserted{(ext)} de $\pi$ il est sous-entendu que \deleted{\illegible{1}} non seulement ils commutent à $\Gamma$, mais encore qu'ils sont dans $N(\pi)$. \inserted{(ce qui implique qu'ils sont dans \unsure{$\widehat{\mathcal T}(\pi)$}, car le centralisateur dans $\Pi_\pi$ de \unsure{tt} sous-groupe ouvert de $\Pi_\pi$ est $(1)$ !)} Dans le cas de $U_{0,3}$, on a $N_{0,3}\simeq\unsure{\mathfrak S_3}\times\Pi_{0,3}$, et on sait que le centre de $\Pi_{03}\simeq\Pi$ est trivial -- donc on trouve bien que le groupe des autom. de cette structure est réduit : \unsure{$\mathfrak S_3$} !

Ce point de vue est néanmoins probablement superflu -- car on \unsure{présume} que pour l'action extérieure donnée de $\Gamma$, il y a une unique \unsure{authentisation} invariante par $\Gamma$, et que les \unsure{hom} \inserted{``admissibles''} des $\Gamma$-groupes ext. ``admissibles'', tels qu'ils ont été définis précédemment, respectent \unsure{automatiquement} cette \unsure{authentisation}. S'il n'en était rien, il faudrait bien entendu introduire les \unsure{authentisations} dans la structure.

\deleted{\unsure{Ici}, \deleted{\unsure{comme} $\Pi_\pi\unsure{=}$} \inserted{si $\Gamma=\unsure{\Pi_{\mathbb Q}}$,} l'\unsure{authentisation} de \deleted{\illegible{1}} $\pi$ \deleted{\unsure{est}} \inserted{doit être} \unsure{décrite} en fait par une prédiscrétification orientée bien déterminée $\alpha$, \deleted{\unsure{du fait qu'}\illegible{2}} \inserted{En fait, \unsure{on va|on se} décrire une prédiscrétification \unsure{orient.} \uline{stricte}} \deleted{\illegible{3} \unsure{point de base}} bien déterminée de $\pi$, Partons d'un relèvement admissible de l'op. ext. donnée}

\margin{485}

\scanpage{487}
\mspage{247}

Je \unsure{me} pose la question de trouver une description plus \deleted{\unsure{simple et}} \inserted{\unsure{simpliste|simple}} des actions extérieures d'un $\Gamma$ sur un $\pi$ qui sont ``admissibles''. \deleted{\illegible{1}} Ici, on va partir d'un $\pi$ dont on \unsure{fixe} déjà une structure à lacets \inserted{et une \unsure{authent.}} (\unsure{invariées|invariés} par $\Gamma$), de \deleted{\illegible{2}} sorte que $\Gamma\subset N(\pi)$, $\Gamma\longrightarrow\Pi_\pi$ \unsure{injective} à image ouverte, \deleted{\illegible{1}} i.e. $\Gamma$ correspond à un scindage partiel (ou germe de scindage) de \deleted{\illegible{1}}
\[
1\longrightarrow\widehat{\mathcal T}(\pi)\longrightarrow N(\pi)\longrightarrow\Pi_\pi\longrightarrow1
\]

\unsure{Soit} $(g,\nu)$ le type de $\pi$; \deleted{\illegible{2}} si $g\geq1$ le critère le plus simple, c'est que les ``\inserted{$\Gamma^{\unsure{\natural}}$-}points'' du $\pi'$ \unsure{déduit} de $\pi$ en bouchant tous les trous, sont distincts. \deleted{\illegible{1} $g\geq1$} (NB si $g=1$, \deleted{\unsure{cette}} \deleted{\unsure{bouchant}} \inserted{en bouchant} tous les trous on tombe dans un cas abélien -- mais ça ne fait rien). Si $g=0$, \deleted{le critère est que} il n'y a pas de condition pour $\unsure{\alpha|\nu}=3$, et si $\nu>3$, \unsure{mettons} à part une \deleted{\illegible{1}} partie $I'\subset I(\pi)$ avec $\mathrm{card}\,I'=3$, la condition c'est que en bouchant les trous \unsure{de|dans} $\unsure{I\setminus I'}$, les \inserted{$\Gamma^{\unsure{\natural}}$-}pts \unsure{du|de} $\pi$ déduits des pts de $\unsure{I\setminus I'}$ sont distincts.

On notera que dans cette optique conjecturale, dans les cas limites $(g,0)$ $(g\geq2)$, $(1,1)$, $(0,3)$, on n'impose aucune condition \inserted{a priori} sur les \unsure{relèvements}.

\margin{486}

\scanpage{488}

C'est peut-être très \unsure{brutal} -- et il se pourrait qu'on trouve des relèvements qui ne \unsure{correspondent} \unsure{pas} à une courbe alg. -- i.e. qui en fait ne sont \unsure{pas} admissibles, même s'ils le paraissent.

Pour y comprendre quelque chose, il me semble qu'il faut revenir à une interprétation des \inserted{\unsure{germes de}} scindages dits ``admissibles'' d'une extension telle que
\diagram{\begin{tikzcd}[ampersand replacement=\&]
1 \arrow[r] \& \widehat{\mathcal T}_{g,\nu} \arrow[r] \& N_{g,\nu} \arrow[r] \& \Pi_{g,\nu} \arrow[r] \& 1 \\
\& \& \& \unsure{\Pi_{\mathbb Q}}\simeq\Pi_{0,3} \arrow[u,"\unsure{\wr}"'] \&
\end{tikzcd}}
comme correspondants aux points algébriques d'une \uline{variété} (\unsure{plutôt ici}, une multiplicité) modulaire \unsure{$M_{g,\nu\,\mathbb Q}$|$M_{g,\nu,\mathbb Q}$}, dont $N_{g,\nu}$ est le groupe fondamental arithmétique, et $\widehat{\mathcal T}_{g,\nu}$ le groupe fondamental géométrique.

J'aimerais examiner cette situation de plus près, par la suite. Pour le moment, il semble prudent de ne pas faire des conjectures hâtives pour \deleted{\illegible{1}} une description \uline{directe} (i.e. \uline{pas} via $U_{03}$) des opérations extérieures ``admissibles''. On \unsure{travaillera} donc pour le moment sur cette notion \unsure{sous la|sur la} \unsure{forme} ``\unsure{constructive}'' (via $U_{03}$).

\margin{487}

\scanpage{489}
\mspage{248}

Si on a un $\pi$ \inserted{ext.} de type \unsure{$(g\nu)$|$(g,\nu)$} avec op. extérieure admissible de $\Gamma\subset\Pi$, on prévoit qu'il y aura une prédiscrétification \unsure{orientée} canonique sur \unsure{$\pi_x$} (pas \unsure{invariant|invariante} par $\Gamma$, bien sûr \unsure{!|)} -- mais telle que l'\unsure{authentisation} correspondante le soit). On \unsure{va} \inserted{$\overline{\mathrm m}$}, pour une action \inserted{effective} admissible de \unsure{$\Gamma^{\natural}$|$\Gamma^{\flat}$} sur $\pi$, \inserted{(relevant l'action ext. donnée)}, définir une discrétification \inserted{or.} \unsure{correspondante} canonique $\pi_0\subset\pi$ (\unsure{remplacée} par une \unsure{conjuguée}, quand on conjugue \deleted{\unsure{l'action}} le relèvement $\unsure{\Gamma^{\natural}|\Gamma^{\flat}}\to\mathcal E$ par un $g\in\pi$) -- toutes ces discrétifications or. (correspondantes aux différents ``pts'' de $B_{\pi,\unsure{\Gamma^{\natural}|\Gamma^{\flat}}}$) définissent une \deleted{$\overline{\mathrm m}$} \inserted{pré-}discrétification orientée -- et $\overline{\mathrm m}$ une $\overline{\mathrm m}$ prédiscr. or. \uline{stricte}.

L'action effective de \unsure{$\Gamma^{\natural}$|$\Gamma^{\flat}$} sur $\pi$ peut s'obtenir (à isom. près) \inserted{comme suit:} en prenant un sous-groupe \inserted{ouvert} $\mathcal E'\subset\mathcal E_{03}$, d'où extension
\[
1\to\pi'\to\mathcal E'\to\Gamma'\to1
\]
avec $\pi'\subset\widehat{\unsure{\pi|\Pi}}_{03}$, donc $\pi'=\widehat{\pi}'_0$ $(\pi'_0=\pi'\cap\unsure{\pi|\Pi}_{03})$. On a une opération de $\Gamma'$ sur $I(\pi')=I(\pi'_0)$, on la prend triviale (quitte à passer à un ss-groupe ouvert de $\Gamma'$), on choisit $i\in I'\subset I(\pi')$, on bouche les trous en $I'$, \deleted{et on regarde} \inserted{d'où $\pi$, une \deleted{action ext.}} discrétification \inserted{orientée} $\pi_0$, et une action ext. de $\Gamma'$ dessus, qui est relevée en une action effective grâce à $i$. On trouve bien ainsi une

\margin{488}

\scanpage{490}

``réseau'' \inserted{-- une discrétification or.} dans $\pi$ -- je dis qu'elle ne dépend pas des choix qui \unsure{ont} été faits -- en particulier, que les \unsure{autom.|autos} de $(\pi,a,\unsure{\Gamma^{\natural}|\Gamma^{\flat}})$ \unsure{transforment} $\pi_0$ en lui-$\overline{\mathrm m}$. De plus, si on a deux pts \inserted{$x,y$} de $B_{\pi,\unsure{\Gamma^{\natural}|\Gamma^{\flat}}}$, d'où $\pi(x)$, $\pi(y)$, parmi les \unsure{classes} de chemins de $x$ à $y$ (qui \unsure{font|forment} un \unsure{bitorseur} sous $\pi(y),\pi(x)$) il y en a pour lesquels \deleted{\illegible{2}} $\pi(x)\to\pi(y)$ envoie $\pi_0(x)$ dans $\pi_0(y)$, -- quand on les \unsure{prend en masse}, on trouve un sous-groupoïde du groupoïde fondamental $\Pi_1(\pi,\unsure{\Gamma^{\natural}|\Gamma^{\flat}})$, qui est cette fois ci un groupoïde \uline{\unsure{connexe}} à lacets. On \unsure{fera} attention que $\Gamma$ \inserted{n'}opère \uline{pas} sur ce groupoïde, bien qu'il opère sur l'ens. de ses objets.

\margin{489}
